\section{Latent Semantic Indexing}

LSI represents documents and queries as vectors in cartesian space, and measuring the similarity between a query and a document, as the cosine of the angle between their vector representation

literal matching scheme such as VSM suffer from synonyms and noise in documents. LSI overcomes these problems by using statistically derived conceptual indices instead of terms for retrieval.

It uses Singular Value Decomposition (SVD) to transform a high-dimensional term vector (computed from VSM) into a lower-dimensional semantic vector, by projecting the former into a semantic subspace. Each element of a semantic vector corresponds to the importance of an abstract concept in the document or query

Suppose the rank of A is $r$. SVD decomposes $A$ into the product of three matrices:
\begin{subequations}
\begin{equation}
\bvec{A_{m \times n}} = 
\bvec{U_{m \times r}}
\underbrace{
\left(
 \begin{array}{ccc}
 \sigma_1 & \cdots & 0        \\
 \vdots   & \ddots & \vdots   \\
        0 & \cdots & \sigma_r \\
\end{array}
 \right)}_{\Sigma_{r \times r}}
\bvec{V_{r \times n}^T}
\end{equation}
\begin{equation}
\bvec{U}^T \bvec{U} = \bvec{V}^T \bvec{V} = \bvec{I}
\end{equation}
\begin{equation}
\sigma_1 \geq \cdots \geq \sigma_r
\end{equation}
\end{subequations}

LSI approximates the matrix A of rank $r$ with a matrix $A_k$ of lower rank $k$ by omitting \ldots $k$ singular values
\begin{equation}
\bvec{A_k} = \bvec{U_k} \bvec{\Sigma_k} {\bvec{V_k}}^T
\end{equation}

The rows of $\bvec{V_k}\bvec{\Sigma_k}$ are the semantic vectors for documents in the corpus.

By choosing an appropriate $k$ for $A_k$ the important structure of the corpus is retained while the noise or variability in word usage (small $\sigma_i$) is eliminated. Former studies on LSI suggest setting $k$ to a value between $50$ and $350$

LSI is capable of bringing together documents that are semantically related even if they do not share terms, by learning from co-occurring word usage

the $i^{th}$ row of $\bvec{U}$ is a refined representation of term $i$ and the $j^{th}$ row of $\bvec{V}$ is a refined representation of document $j$. Both representations are vectors in a $r$-dimensional subspace.

By choosing $k=2,3$ we can plot and visualize those vectors

LSI maps similar words into similar vectors exploit correlations between terms

\begin{figure}
\begin{center}
\includegraphics[width=130mm]{img/clustering/title-leo2d.pdf}
\caption{2-dimensional representation of documents vectors (rows of \bvec{V}) }%\label{fig:pdfParser_statechart}
\end{center}
\end{figure}

\begin{figure}
\begin{center}
\includegraphics[width=130mm]{img/clustering/term-term_comparison.pdf}
\caption{2-dimensional representation of term vectors (rows of \bvec{U}) }%\label{fig:pdfParser_statechart}
\end{center}
\end{figure}







\begin{figure}
\begin{center}
\includegraphics[width=120mm]{img/clustering/abstract-comparison.pdf}
\caption{abstract-comparison.pdf}%\label{fig:pdfParser_statechart}
\end{center}
\end{figure}

\begin{figure}
\begin{center}
\includegraphics[width=120mm]{img/clustering/abstract-comparison3d.pdf}
\caption{abstract-comparison3d}%\label{fig:pdfParser_statechart}
\end{center}
\end{figure}






\begin{figure}
\begin{center}
\includegraphics[width=120mm]{img/clustering/title-comparison.pdf}
\caption{img/clustering/title-comparison.pdf}%\label{fig:pdfParser_statechart}
\end{center}
\end{figure}






%\begin{figure}
%\begin{center}
%\includegraphics[width=120mm]{img/clustering/title-leo3d.pdf}
%\caption{title-leo3d.pdf}%\label{fig:pdfParser_statechart}
%\end{center}
%\end{figure}

\subsection{Clustering}
\paragraph{Centroid}
\begin{equation}
\mu = \frac{1}{|\omega|} \sum_{\bvec{x} \in \omega}\bvec{x}
\end{equation}

\paragraph{RSS} Residual sum of squares
\begin{subequations}
\begin{equation}
\mbox{RSS}_k = \sum_{\bvec{x} \in \omega_k} | \bvec{x} - \mu(\omega_k) |^2
\end{equation}
\begin{equation}
\mbox{RSS} = \sum_k \mbox{RSS}_k
\end{equation}
\end{subequations}
